
\renewcommand{\arraystretch}{1.5} % Default value: 1

%%% Write here
%%%%%%%%%%%%%%%%%%%%%%%%%%%%%%%%%%%%%
\section{Introduction}
Optical polarimetry, ever since its discovery, has played important roles in gaining fundamental understanding of the nature of electromagnetic waves and answering some of the key questions related to the physics of light. Traditional polarimetry has also long being pursued for numerous practical applications in various branches of science and technology. Quantification of protein properties in solutions, testing of chiral (asymmetric) purity of pharmaceutical drugs, remote sensing in meteorology and astronomy, optical stress analysis of structures, and crystallography of biochemical complexes are just a smattering of its diversified uses. This particular polarimetry experiment involves determination of two important polarization characteristics of medium, namely, the linear retardance and optical rotation (defined subsequently), which has many practical applications.
%%%%%%%%%%%%%%%%%%%%%%%%%%%%%%%%%%%%%
\section{Objective}
The main opjective of this experiment is to determine of linear retardance and Optical rotation of a Spatial Light Modulator (SLM) measuring various Stokes vector elements
%%%%%%%%%%%%%%%%%%%%%%%%%%%%%%%%%%%%%
\section{Equipment Required}
The following equipments are required for this experiment
\begin{itemize}
    \item Diode laser (650 nm)
    \item Spatial Light Modulator (SLM)
    \item Two Polarizer with mounts
    \item Photo-detector
    \item Optical table and mounts
    \item Computer with software for data acquisition and analysis
    \item Power supply for the SLM
    \item Glass plate with stage having rotational circular scale(in degree)
\end{itemize}
%%%%%%%%%%%%%%%%%%%%%%%%%%%%%%%%%%%%%
\section{Theory}


\subsection{Basics of Polarization}
Christian Huygens was the first to suggest that light was not a scalar quantity based on his work on the propagation of light through crystals; it appeared that light had `sides' in the words of Newton. This vectorial nature of light is called polarization \cite{Goldstein2003PolarizedLight}.

In a Cartesian system the components $u_x(\mathbf{r},t)$ and $u_y(\mathbf{r},t)$ are said to be the transverse components, and the component $u_z(\mathbf{r},t)$  is said to be the longitudinal component when the propagation is in the z direction. Thus we can write

\begin{subequations}
\begin{align}
u_x(\mathbf{r},t) &= u_{0x}\cos(\omega t - \mathbf{k}\cdot\mathbf{r} + \delta_x), \label{eq:u_x}\\
u_y(\mathbf{r},t) &= u_{0y}\cos(\omega t - \mathbf{k}\cdot\mathbf{r} + \delta_y), \label{eq:u_y}\\
u_z(\mathbf{r},t) &= u_{0z}\cos(\omega t - \mathbf{k}\cdot\mathbf{r} + \delta_z). \label{eq:u_z}
\end{align}
\end{subequations}

Experimentally found that light consisted only of the transverse components \cref{eq:u_x} and \cref{eq:u_y}. If we take the direction of propagation to be in the z direction, then the optical field in free space must be described only by \cref{eq:u_x} and \cref{eq:u_y}. where $u_{0x}$ and $u_{0y}$ are the maximum amplitudes and $\delta_x$ and $\delta_y$ are arbitrary phases. These two \cref{eq:u_x} and \cref{eq:u_y} are are spoken of as the polarized or polarization components of the optical field. Now, the transverse components are represented by

\begin{subequations}\label{eq:E_trans}
\begin{align}
E_x(z,t) &= E_{0x}\cos(\tau+\delta_x), \label{eq:E_x}\\
E_y(z,t) &= E_{0y}\cos(\tau+\delta_y), \label{eq:E_y}
\end{align}
\end{subequations}

where $\tau=\omega t-kz$ is the propagator. The subscripts $x$ and $y$ refer to the components in the $x$ and $y$ directions. $E_{0x}$ and $E_{0y}$ are the maximum amplitudes, and $\delta_x$ and $\delta_y$ are the phases, respectively. As the field propagates, $E_x(z,t)$ and $E_y(z,t)$ give rise
to a resultant vector. This vector describes a locus of points in space, and the curve generated by those points will now be derived. In order to do this \cref{eq:E_x} and \cref{eq:E_y} are written as

\begin{subequations}
\begin{align}
\frac{E_x}{E_{0x}} &= \cos\tau\cos\delta_x- \sin\tau\sin\delta_x, \\
\frac{E_y}{E_{0y}} &= \cos\tau\cos\delta_y - \sin\tau\sin\delta_y. 
\end{align}
\end{subequations}


\begin{figure}[H]
    \centering
    \begin{subfigure}[t]{0.45\textwidth}
        \centering
        \includegraphics[width = \textwidth]{Image/polarization2.jpg}
        \caption{An elliptically polarized wave and the polarization ellipse}
        \label{fig:polarization_ellipse}
    \end{subfigure}
    \hfill
    \begin{subfigure}[t]{0.45\textwidth}
        \centering
        \includegraphics[width = \textwidth]{Image/polarization1.jpg}
        \caption{Propagation of the transverse optical field}
        \label{fig:propagation_wave}
    \end{subfigure}
    \caption{Polarization of light}  
    % \label{fig:enter-label}
\end{figure}



Hence,

\begin{subequations}\label{eq:polarization-relations}
\begin{align}
\frac{E_x}{E_{0x}}\sin\delta_y
-\frac{E_y}{E_{0y}}\sin\delta_x
&= \cos\tau\,\sin(\delta_y-\delta_x),
\label{eq:cos-relation}\\
\frac{E_x}{E_{0x}}\cos\delta_y
-\frac{E_y}{E_{0y}}\cos\delta_x
&= \sin\tau\,\sin(\delta_y-\delta_x).
\label{eq:sin-relation}
\end{align}
\end{subequations}

Squaring \eqref{eq:cos-relation} and \eqref{eq:sin-relation} and adding gives

\begin{subequations}\label{eq:polarization-ellipse}
\begin{align}
\frac{E_x^2}{E_{0x}^2}
+\frac{E_y^2}{E_{0y}^2}
-2\frac{E_x}{E_{0x}}\frac{E_y}{E_{0y}}\cos\delta
&=\sin^2\delta,
\label{eq:ellipse-equation}\\
\delta &= \delta_y-\delta_x.
\label{eq:phase-difference}
\end{align}
\end{subequations}

\cref{eq:ellipse-equation} is recognized as the equation of an ellipse and shows that at any instant of time the locus of points described by the optical field as it propagates is an ellipse (see \cref{fig:polarization_ellipse}). This behavior is spoken of as optical polarization, and \cref{eq:ellipse-equation} is called the \emph{polarization ellipse} \cite{Goldstein2003PolarizedLight}.

Also, we can write some more parameters of ellipse in terms of $\delta$ and $\alpha$ where $\tan{\alpha} = E_{0y}/E_{0x}$

\begin{subequations}
    \begin{align}
    \tan 2\psi &= (\tan 2\alpha)\cos\delta, 
    \qquad 0 \le \psi \le \pi, \\
    \sin 2\chi &= (\sin 2\alpha)\sin\delta,
    \qquad -\frac{\pi}{4} \le \chi \le \frac{\pi}{4}.
    \end{align}
\end{subequations}

\subsection{Stokes Vector and Muller Matrix}

The \emph{Stokes vector} is a convenient method for describing the polarization state of light using measurable intensity quantities. It consists of four parameters $(S_0, S_1, S_2, S_3)$, where $S_0$ represents the total intensity, while $S_1$, $S_2$, and $S_3$ characterize the linear and circular polarization components. Since these parameters can be obtained experimentally using polarizers and wave plates, the Stokes formalism is widely used in polarization analysis and optical characterization.

We now write the quantities inside the parentheses as
\begin{subequations}
    \begin{align}
    S_0 &= E_{0x}^{\,2} + E_{0y}^{\,2}, \label{eq:S_0}\\
    S_1 &= E_{0x}^{\,2} - E_{0y}^{\,2}, \label{eq:S_1}\\
    S_2 &= 2E_{0x}E_{0y}\cos\delta, \label{eq:S_2}\\
    S_3 &= 2E_{0x}E_{0y}\sin\delta, \label{eq:S_3}
    \end{align}
\end{subequations}
and then using the \cref{eq:polarization-ellipse} we can write
\begin{equation}
S_0^{2} = S_1^{2} + S_2^{2} + S_3^{2}.
\end{equation}
Also the other parameters of the polarization ellipse can be write interms of Stokes vectors
\begin{subequations}
    \begin{align}
        \tan{2\Psi} = \frac{S_{2}}{S_{1}}, \\
        \sin{2\chi} = \frac{S_{3}}{S_{0}}.
    \end{align}
\end{subequations}
And the degree of polarization are defined as following

\begin{equation}
P = \frac{I_{\mathrm{pol}}}{I_{\mathrm{tot}}}
= \frac{\left(S_1^{2}+S_2^{2}+S_3^{2}\right)^{1/2}}{S_0},
\qquad 0 \le P \le 1 .
\end{equation}

where $I_{\mathrm{pol}}$ is the intensity of the sum of the polarization components and $I_{\mathrm{tot}}$ is the total intensity of the beam. The value of $P=1$ corresponds to completely polarized light, $P=0$ corresponds to unpolarized light, and $0<P<1$ corresponds to partially polarized light.

Thus these parameters can be directly determined by the following six intensity measurements $(I)$ performed with ideal polarizers: $I_H$, horizontal linear polarizer $(0^\circ)$; $I_V$, vertical linear polarizer $(90^\circ)$; $I_P$, $45^\circ$ linear polarizer; $I_M$, $135^\circ$ ($-45^\circ$) linear polarizer; $I_R$, right circular polarizer; and $I_L$, left
circular polarizer.

\begin{equation}
\mathbf{S} =
\begin{bmatrix}
I \\ Q \\ U \\ V
\end{bmatrix}
=
\begin{bmatrix}
I_H + I_V \\
I_H - I_V \\
I_P - I_M \\
I_R - I_L
\end{bmatrix}.
\end{equation}

\begin{table}[H]
\centering
\caption{Stokes vectors for different polarization states}
\begin{tabular}{|c|c|c|c|c|c|c|c|}
\toprule
State & H & V & P & M & R & L & Elliptical \\
\midrule

$\mathbf{S}$ &
$\begin{pmatrix}1\\[2pt]1\\[2pt]0\\[2pt]0\end{pmatrix}$ &
$\begin{pmatrix}1\\[2pt]-1\\[2pt]0\\[2pt]0\end{pmatrix}$ &
$\begin{pmatrix}1\\[2pt]0\\[2pt]1\\[2pt]0\end{pmatrix}$ &
$\begin{pmatrix}1\\[2pt]0\\[2pt]-1\\[2pt]0\end{pmatrix}$ &
$\begin{pmatrix}1\\[2pt]0\\[2pt]0\\[2pt]1\end{pmatrix}$ &
$\begin{pmatrix}1\\[2pt]0\\[2pt]0\\[2pt]-1\end{pmatrix}$ &
$\begin{pmatrix}
1\\[2pt]
\cos 2\alpha\\[2pt]
\sin 2\alpha \cos\delta\\[2pt]
\sin 2\alpha \sin\delta
\end{pmatrix}$ \\

\bottomrule
\end{tabular}

\label{tab:stokes_vectors}
\end{table}

The \emph{Mueller matrix} provides a general framework for describing how an optical system modifies the polarization state of light. It relates the incident and transmitted Stokes vectors through a $4\times4$ real matrix and can represent fully polarized, partially polarized, and unpolarized light. Owing to its basis in measurable intensity quantities, the Mueller matrix formalism is widely used for polarization analysis and characterization of optical elements.

In matrix form this can be write as 
\begin{equation}
\begin{pmatrix}
S_0' \\
S_1' \\
S_2' \\
S_3'
\end{pmatrix}
=
\begin{pmatrix}
m_{00} & m_{01} & m_{02} & m_{03} \\
m_{10} & m_{11} & m_{12} & m_{13} \\
m_{20} & m_{21} & m_{22} & m_{23} \\
m_{30} & m_{31} & m_{32} & m_{33}
\end{pmatrix}
\begin{pmatrix}
S_0 \\
S_1 \\
S_2 \\
S_3
\end{pmatrix}
\end{equation}

or
\begin{equation}
\mathbf{S}' = \mathbf{M}\cdot \mathbf{S}.
\end{equation}

\subsection{Medium Polarization Characteristics}

A transparent optical medium is generally characterized by three
fundamental polarization properties: diattenuation, depolarization, and retardance. In this experiment, the primary focus is on determining the retardance of the material. Retardance originates from the birefringent nature of a medium and produces a phase difference between orthogonal polarization components of light. This phase difference depends on the thickness and optical properties of the material. Retardance can be classified into two main types.

\begin{enumerate}
\item \textbf{Linear retardance:}
Linear retardance arises from a phase delay between two \textit{orthogonal linear polarization components}. It results from linear birefringence, which occurs due to anisotropy in the medium, leading to different refractive indices for orthogonal polarization directions. For a medium of thickness $L$, the linear retardance is given by
\begin{equation}
\delta = \frac{2\pi}{\lambda}(n_x - n_y)L,
\end{equation}

where $\lambda$ denotes the wavelength of light, and $n_x$ and $n_y$ represent the refractive indices along two orthogonal axes.

\item \textbf{Circular retardance:}
Circular retardance corresponds to \textit{the phase difference between right- and left-circularly polarized light}. This effect originates from circular birefringence, typically present in optically active media lacking mirror symmetry. Such anisotropy produces different refractive indices for the two circular polarization states. For a propagation length $L$, the circular retardance is expressed as

\begin{equation}
\delta_c = \frac{2\pi}{\lambda}(n_R - n_L)L,
\end{equation}

where $n_R$ and $n_L$ are the refractive indices for right- and
left-circularly polarized light, respectively. \textit{Circular retardance manifests as optical rotation}, corresponding to the rotation of the plane of linear polarization about the propagation axis, with rotation angle

\begin{equation}
\psi = \frac{\delta_c}{2}.
\end{equation}

\end{enumerate}

\subsection{Spatial Light Modulator (SLM)}
The \emph{Spatial Light Modulator} (SLM) is an opto-electronic device used to control the spatial properties of a light beam by modulating its amplitude, phase, or polarization. Most commonly, SLMs are based on liquid crystal technology, where a pixelated layer of liquid crystals is placed between polarizing elements. When an external voltage is applied to each pixel, the orientation of the liquid crystal molecules changes, altering the polarization state of the transmitted light and thereby controlling its intensity or phase. Because each pixel can be addressed electronically through a computer, SLMs enable dynamic and real-time manipulation of optical wavefronts, allowing the 
generation of programmable amplitude and phase patterns for various applications in beam shaping, holography, and optical experiments.

Since the SLM is made of liquid crystals which rotate and align according to the voltage applied to them, the general form of the Mueller matrix of an SLM can be written as a retarder matrix multiplied with a rotation matrix. Now the Muller matrix for a retarder (or compensator) with fast axis toward x-axis can be written as 

\begin{equation}
M_c =
\begin{pmatrix}
1 & 0 & 0 & 0 \\
0 & 1 & 0 & 0 \\
0 & 0 & \cos\delta & \sin\delta \\
0 & 0 & -\sin\delta & \cos\delta
\end{pmatrix}
\end{equation}

For $\theta$ angle rotated wave plate we can write 

\begin{subequations}
    \begin{align}
        M_c(\delta,2\theta) &= M_R(-2\theta)M_cM_R(2\theta) \\ 
        M_R(2\theta) &=
        \begin{pmatrix}
        1 & 0 & 0 & 0 \\
        0 & \cos 2\theta & \sin 2\theta & 0 \\
        0 & -\sin 2\theta & \cos 2\theta & 0 \\
        0 & 0 & 0 & 1
        \end{pmatrix}
    \end{align}
\end{subequations}

So we can write 

\begin{equation}
    M_c(\delta,2\theta)=
    \begin{pmatrix}
    1 & 0 & 0 & 0 \\
    0 & \cos^2 2\theta + \cos\delta\,\sin^2 2\theta &
    (1-\cos\delta)\sin 2\theta \cos 2\theta &
    -\sin\delta \sin 2\theta \\
    0 & (1-\cos\delta)\sin 2\theta \cos 2\theta &
    \sin^2 2\theta + \cos\delta\,\cos^2 2\theta &
    \sin\delta \cos 2\theta \\
    0 & \sin\delta \sin 2\theta &
    -\sin\delta \cos 2\theta &
    \cos\delta
    \end{pmatrix}.
\end{equation}

Now at the end for SLM we can write 

\begin{subequations}
    \begin{align}
        & M_{\mathrm{SLM}} = M_{c}(\delta,2\theta) \times M_{R}(-2\Psi) \\
        & = 
        \begin{pmatrix}
        1 & 0 & 0 & 0 \\
        0 & \cos^{2}2\theta + \sin^{2}2\theta \cos\delta
        & \sin2\theta\cos2\theta(1-\cos\delta)
        & -\sin2\theta\sin\delta \\
        0 & \sin2\theta\cos2\theta(1-\cos\delta)
        & \sin^{2}2\theta + \cos^{2}2\theta \cos\delta
        & \cos2\theta\sin\delta \\
        0 & \sin2\theta\sin\delta
        & -\cos2\theta\sin\delta
        & \cos\delta
        \end{pmatrix}
        \begin{pmatrix}
        1 & 0 & 0 & 0 \\
        0 & \cos2\psi & -\sin2\psi & 0 \\
        0 & \sin2\psi & \cos2\psi & 0 \\
        0 & 0 & 0 & 1
        \end{pmatrix} \\
        &\begin{pmatrix}
        1 & 0 & 0 & 0 \\
        0 & \cos 2\theta \cos(2\theta-2\psi) + \sin 2\theta \cos\delta \sin(2\theta-2\psi) &  -\cos\delta \sin 2\theta \cos(2\theta-2\psi) + \cos 2\theta \sin(2\theta-2\psi)
        &  -\sin\delta \sin 2\theta \\
        0 & \sin 2\theta \cos(2\theta-2\psi)- \cos 2\theta \cos\delta \sin(2\theta-2\psi) & \cos\delta \cos 2\theta \cos(2\theta-2\psi) + \sin 2\theta \sin(2\theta-2\psi) & \cos 2\theta \sin\delta \\
        0 &  \sin\delta \sin(2\theta-2\psi) & -\cos(2\theta-2\psi)\sin\delta & \cos\delta
        \end{pmatrix}
    \end{align}
\end{subequations}

Now as per the experiment once we take $I_{H}$ polarized light and another time we take $I_P$ polarized light. By applying the SLM matrix we can find

\begin{equation}
    S_H = M_{\mathrm{SLM}} \times \begin{pmatrix}1\\[2pt]1\\[2pt]0\\[2pt]0\end{pmatrix} = \begin{pmatrix}
    1 \\
    \cos 2\theta \cos(2\theta-2\psi)
    + \sin 2\theta \cos\delta \sin(2\theta-2\psi) \\
    \sin 2\theta \cos(2\theta-2\psi)
    - \cos 2\theta \cos\delta \sin(2\theta-2\psi) \\
    \sin(2\theta-2\psi)\sin\delta
    \end{pmatrix}
\end{equation}

and, 

\begin{equation}
        S_P = M_{\mathrm{SLM}} \times \begin{pmatrix}1\\[2pt]0\\[2pt]1\\[2pt]0\end{pmatrix} = \begin{pmatrix}
        1 \\
        -\cos\delta \sin 2\theta \cos(2\theta-2\psi)
        + \cos 2\theta \sin(2\theta-2\psi) \\
        \cos\delta \cos 2\theta \cos(2\theta-2\psi)
        + \sin 2\theta \sin(2\theta-2\psi) \\
        -\cos(2\theta-2\psi)\sin\delta
        \end{pmatrix}
\end{equation}

So finally we can write 

\begin{subequations}
    \begin{align}
    S_H(2) + S_P(3) &= (1+\cos\delta)\cos 2\psi = X_1, \\
    S_H(3) - S_P(2) &= (1+\cos\delta)\sin 2\psi = Y_1,
    \end{align}
\end{subequations}

\begin{equation}
\boxed{\delta = \cos^{-1}\!\left[\sqrt{X_1^{2}+Y_1^{2}} - 1 \right]}
\end{equation}

\begin{equation}
\boxed{\psi = \frac{1}{2}\tan^{-1}\!\left(\frac{Y_1}{X_1}\right)}.
\end{equation}

where $\delta$ is \emph{linear retardance} and $\Psi$ is \emph{optical rotation}.
%%%%%%%%%%%%%%%%%%%%%%%%%%%%%%%%%%%%%
\section{Set up and Procedure}
This experiment has two parts. In the first part, we did measure the angles of polarizers (P1 and P2) for Horizontal polarization. And in the later part we did Stokes polarimetry for SLM. 
\subsection*{Part 1}


\begin{figure}[H]
    \centering
    \includegraphics[width = 0.6\textwidth]{Image/part1.JPG}
    \caption{Setup to find the axis of the first polarizer (P1) using Brewster's angle approch}
\end{figure}

\begin{figure}[H]
    \centering
    \includegraphics[width = \textwidth]{Image/schematic.jpg}
    \caption{Schematic of the final setup}
\end{figure}

\begin{figure}[H]
    \centering
    \includegraphics[width = 0.8\textwidth]{Image/part2.JPG}
    \caption{Final setup to find the Stokes measurement}
\end{figure}

%%%%%%%%%%%%%%%%%%%%%%%%%%%%%%%%%%%%%
\section{Data analysis and calculation} 
%%%%%%%%%%%%%%%%%%%%%%%%%%%%%%%%%%%%%
\section{Conclusion}
%%%%%%%%%%%%%%%%%%%%%%%%%%%%%%%%%%%%%